\chapter{数値出力の記録}\label{app:outputs}
本文の数値の出典として、各シミュレーションの標準出力（\texttt{data/*\_summary.txt}）をそのまま収録する。Sim01・Sim02 の数値は第~\ref{ch:planar}章の本文に記載した。

\section{Sim03（\texttt{data/sim03\_summary.txt}）}
{\footnotesize\raggedright\sloppy\setlength{\parindent}{0pt}\setlength{\parskip}{0pt}\emergencystretch=3em
=== Part A ===\par
Sim02 の長方形ループ: +2.98°/周\par
最適円ループ: 中心(腕 -20°, 脚 -40°), 半径 70° → +6.70°/周\par
可動域内の理論上界（同符号曲率の全領域）: +9.57° / -2.23°\par
RW の Γ\_w = I\_w/A = 0.00140 (一定) → ホイール1回転あたり 0.51°\par
\mbox{}\par
=== Part B: 90° rest-to-rest ===\par
RW   : ホイール 178.0 回転, 最大 2003 rpm, 所要 10 s → 90.00°\par
AMBAC: 最適ループ 13 周, 所要 26 s → 87.16°（整数周回による端数あり）, 途中の最大姿勢 113.2°\par
\mbox{}\par
=== Part C: 外乱トルク 0.010 N m 下の姿勢保持 ===\par
RW    の角運動量容量 h = I\_w ω\_max          = 8.640 kg m²/s → 保持限界 864 s\par
AMBAC 保持中: 周回平均姿勢 |$\bar\theta$| ≤ 0.30°, 瞬時姿勢の揺れ ±23.3°\par
AMBAC の角運動量容量 h = f\_max |Δθc| / <1/A> = 1.180 kg m²/s → 保持限界 118 s\par
}

\section{Sim04（\texttt{data/sim04\_summary.txt}）}
{\footnotesize\raggedright\sloppy\setlength{\parindent}{0pt}\setlength{\parskip}{0pt}\emergencystretch=3em
=== Part A: 動力学モデルの検証（人間大, 0.5 Hz, 2 周, トルクのみ入力・開ループ）===\par
関節軌道の追従誤差 max      : 2.63e-08°\par
2 周後の正味回転 (動力学)   : +13.4086°  (運動学 Sim03: +13.4086°)\par
角運動量 |L| max            : 1.95e-09 kg m²/s\par
エネルギー収支 |ΔT - ∫P dt| : 9.28e-04 J\par
\mbox{}\par
=== Part B: 人間大（76 kg）, 最適ループ ===\par
0.5 Hz での必要トルク peak: 肩 8.2 N m, 股 28.8 N m\par
0.5 Hz でのパワー peak: 肩 21.3 W, 股 47.1 W / 平均正仕事率 計 18.5 W\par
トルク上限 (肩 40, 股 120 N m) での f\_max = 1.02 Hz\par
90° 旋回: f\_max で 13.2 s, 0.5 Hz で消費 497 J（回生なし）\par
\mbox{}\par
=== Part C: スケーリング（トルク上限 ∝ s³ と仮定）===\par
 s   質量[t]  股τ@0.5Hz[N m]  f\_max[Hz]  90°旋回[s]  h[kg m²/s]  平均正パワー@f\_max[kW]\par
 0.5    0.009          0.901      2.040        6.6      0.0752       0.0393\par
 1.0    0.076           28.8      1.020       13.2         1.2        0.157\par
 2.0    0.608            922      0.510       26.3        19.3        0.629\par
 5.0    9.500       9.01e+04      0.204       65.8         752         3.93\par
10.0   76.000       2.88e+06      0.102      131.6     1.2e+04         15.7\par
20.0  608.000       9.22e+07      0.051      263.2    1.93e+05         62.9\par
数値で得た指数: τ@固定f ∝ s\^{}5.00, f\_max ∝ s\^{}-1.00, h ∝ s\^{}4.00  (理論: 5, -1, 4)\par
\mbox{}\par
18 m 級 (s=10, 76 t):\par
  人間と同じ 0.5 Hz で回すのに必要な股関節トルク 2.88 MN m（人間の 100000 倍）\par
  相似則トルク（股 120 N m × s³ = 120 kN m）では f\_max 0.102 Hz, 90° 旋回 132 s\par
  90° を 10 s で旋回するには f = 1.34 Hz → 股関節トルク 20.8 MN m（相似則の 173 倍）\par
\mbox{}\par
=== Part D: 18 m 級の 90° 旋回 スラスタ vs AMBAC ===\par
慣性モーメント A ≈ 1.03e+06 kg m², スラスタ腕 5.0 m, Isp 300 s\par
  T =    5.0 s: 推進剤   88.26 kg / 回\par
  T =   10.0 s: 推進剤   44.13 kg / 回\par
  T =   30.0 s: 推進剤   14.71 kg / 回\par
  T =  131.6 s: 推進剤    3.35 kg / 回\par
  AMBAC (T = 132 s): 推進剤 0 kg、関節の正仕事 約 2.07 MJ（回生なし）\par
}

\section{Sim05A（\texttt{data/sim05a\_summary.txt}）}
{\footnotesize\raggedright\sloppy\setlength{\parindent}{0pt}\setlength{\parskip}{0pt}\emergencystretch=3em
[plain] iter    0  reward -4.1578  (0 s)\par
[plain] iter   50  reward -3.0131  (20 s)\par
[plain] iter  100  reward -3.0059  (40 s)\par
[plain] iter  150  reward -3.0044  (60 s)\par
[plain] iter  200  reward -3.0037  (78 s)\par
[plain] iter  249  reward -3.0012  (97 s)\par
[cpg] iter    0  reward -3.9020  (0 s)\par
[cpg] iter   50  reward -3.0323  (20 s)\par
[cpg] iter  100  reward -3.0089  (40 s)\par
[cpg] iter  150  reward -3.0055  (61 s)\par
[cpg] iter  200  reward -3.0034  (80 s)\par
[cpg] iter  249  reward -3.0038  (101 s)\par
[tg] iter    0  reward -4.7013  (0 s)\par
[tg] iter   50  reward -0.5296  (10 s)\par
[tg] iter  100  reward -0.4706  (20 s)\par
[tg] iter  150  reward -0.4694  (30 s)\par
[tg] iter  200  reward -0.4576  (40 s)\par
[tg] iter  249  reward -0.4581  (49 s)\par
\mbox{}\par
=== 評価（25 s, 指標: 終端誤差 / 最後 5 s の平均誤差）===\par
目標   +90°:  plain  61.8° /  61.8°  cpg  61.8° /  61.8°  tg   0.2° /   0.2°  expert   8.7° /   8.7°\par
目標   -90°:  plain  49.7° /  49.7°  cpg  49.7° /  49.7°  tg   0.1° /   0.1°  expert   7.9° /   7.8°\par
目標   +45°:  plain  16.8° /  16.8°  cpg  16.8° /  16.8°  tg   0.1° /   0.1°  expert   6.1° /   6.1°\par
目標  +150°:  plain 121.8° / 121.8°  cpg 121.8° / 121.8°  tg   0.4° /   0.4°  expert  21.3° /  17.3°\par
[plain] 格子モデルと厳密モデルの差 max 0.010°\par
[cpg] 格子モデルと厳密モデルの差 max 0.009°\par
[tg] 格子モデルと厳密モデルの差 max 0.319°\par
学習と同じ評価関数での報酬: plain -3.001, cpg -3.004, tg -0.458, expert -1.177\par
[tg] 学習したループ: 中心 (腕 -11°, 脚 -10°), 振幅 (120°, 128°), δ 1° → 1 周あたり +9.34°（Sim03 最適円 +6.70°, 可動域内上界 +9.57°）\par
[tg] 目標 +90°: 関節軌道の符号付き掃引面積 +33.5 rad²（正=反時計回り）\par
[tg] 目標 -90°: 関節軌道の符号付き掃引面積 -33.3 rad²（正=反時計回り）\par
}

\section{Sim05B（\texttt{data/sim05b\_summary.txt}）}
{\footnotesize\raggedright\sloppy\setlength{\parindent}{0pt}\setlength{\parskip}{0pt}\emergencystretch=3em
18 m 級: Ā = 1.01e+06 kg m², Δθc = 6.70°/周, f\_max = 0.102 Hz\par
  AMBAC 単独の 90° 旋回時間 ≈ 132 s, RW: h\_max 8.64e+04, τ\_max 200 N m\par
  スラスタ: τ\_max 2.5e+05 N m, 推進剤 0.068 g/(N m s)\par
\mbox{}\par
=== Part B: 90° 旋回の最小推進剤 [kg] ===\par
   T[s]       thrusters         thr+RW      thr+AMBAC   thr+AMBAC+RW\par
    5.2           67.09          66.90          61.18          61.09\par
    8.1           30.11          30.04          28.12          28.06\par
   12.5           18.03          17.93          16.19          16.10\par
   19.4           11.38          11.25           9.70           9.57\par
   30.1            7.27           7.06           5.60           5.40\par
   46.7            4.66           4.34           3.00           2.68\par
   72.5            2.99           2.49           1.34           0.85\par
  112.5            1.92           1.16           0.28           0.00\par
  174.5            1.24           0.05           0.00           0.00\par
  270.7            0.80           0.00           0.00           0.00\par
  thrusters     : 推進剤ゼロ不可\par
  thr+RW        : 推進剤ゼロで旋回できる最短時間 ≈ 183 s（格子上）\par
  thr+AMBAC     : 推進剤ゼロで旋回できる最短時間 ≈ 137 s（格子上）\par
  thr+AMBAC+RW  : 推進剤ゼロで旋回できる最短時間 ≈ 97 s（格子上）\par
\mbox{}\par
=== Part C: 学習 ===\par
iter    0  reward  -103.55  姿勢誤差  75.20°  残留角速度 1.348°/s  推進剤   2.78 kg  (1 s)\par
iter  100  reward    -9.99  姿勢誤差   0.07°  残留角速度 0.121°/s  推進剤  14.99 kg  (86 s)\par
iter  200  reward    -5.39  姿勢誤差   1.20°  残留角速度 0.026°/s  推進剤   7.35 kg  (174 s)\par
iter  300  reward    -3.17  姿勢誤差   0.20°  残留角速度 0.009°/s  推進剤   5.56 kg  (259 s)\par
iter  400  reward    -3.40  姿勢誤差   0.62°  残留角速度 0.009°/s  推進剤   5.20 kg  (342 s)\par
iter  500  reward    -2.56  姿勢誤差   0.06°  残留角速度 0.003°/s  推進剤   4.91 kg  (427 s)\par
iter  600  reward    -3.17  姿勢誤差   0.55°  残留角速度 0.006°/s  推進剤   5.00 kg  (513 s)\par
iter  700  reward    -2.44  姿勢誤差   0.00°  残留角速度 0.009°/s  推進剤   4.50 kg  (598 s)\par
iter  799  reward    -2.60  姿勢誤差   0.09°  残留角速度 0.002°/s  推進剤   4.91 kg  (683 s)\par
参考（公称・平均化モデルの最適値, T = 60 s）: スラスタのみ 3.62 kg, スラスタ+AMBAC 1.97 kg\par
\mbox{}\par
=== Part C: 評価（T = 60 s, 指標: 終端姿勢誤差 / 残留角速度 / 推進剤）===\par
[nominal        ] open-loop OC       : 誤差 平均  17.44° (95\%  17.44°)  残留 0.0000°/s  推進剤   1.97 kg  AMBAC  6.1 周  (0.0 s)\par
[nominal        ] MPC                : 誤差 平均  10.54° (95\%  10.54°)  残留 0.0000°/s  推進剤   3.65 kg  AMBAC  5.4 周  (0.2 s)\par
[nominal        ] MPC + wobble comp. : 誤差 平均   0.02° (95\%   0.02°)  残留 0.0000°/s  推進剤   3.05 kg  AMBAC  0.6 周  (0.4 s)\par
[nominal        ] learned (ES)       : 誤差 平均   0.09° (95\%   0.09°)  残留 0.0024°/s  推進剤   4.66 kg  AMBAC  0.3 周  (0.0 s)\par
[randomized(100)] open-loop OC       : 誤差 平均  30.51° (95\%  68.75°)  残留 0.8956°/s  推進剤   1.97 kg  AMBAC  6.1 周  (0.0 s)\par
[randomized(100)] MPC                : 誤差 平均  14.57° (95\%  23.33°)  残留 0.0303°/s  推進剤   5.41 kg  AMBAC  2.9 周  (14.5 s)\par
[randomized(100)] MPC + wobble comp. : 誤差 平均   0.19° (95\%   0.69°)  残留 0.0307°/s  推進剤   5.56 kg  AMBAC  0.8 周  (41.8 s)\par
[randomized(100)] learned (ES)       : 誤差 平均   0.09° (95\%   0.10°)  残留 0.0024°/s  推進剤   4.92 kg  AMBAC  0.3 周  (0.1 s)\par
}

\section{Sim06（\texttt{data/sim06\_summary.txt}）}
{\footnotesize\raggedright\sloppy\setlength{\parindent}{0pt}\setlength{\parskip}{0pt}\emergencystretch=3em
=== Part A: 荷物運び（前方地面 → 頭上 → 後部デッキ）, 総重量 97 kg, 荷物 10 kg ===\par
最小躍度の単純軌道: 転倒しない最短 1.62 s, 滑らない最短 1.09 s（μ = 0.6）\par
  T = 4.00 s: ZMP [-0.122, +0.166] m（接地範囲 ±0.4）, 必要摩擦係数 max 0.04\par
  T = 2.00 s: ZMP [-0.208, +0.271] m（接地範囲 ±0.4）, 必要摩擦係数 max 0.15\par
  T = 1.62 s: ZMP [-0.282, +0.350] m（接地範囲 ±0.4）, 必要摩擦係数 max 0.23\par
  T = 1.20 s: ZMP [-0.491, +0.560] m（接地範囲 ±0.4）, 必要摩擦係数 max 0.46\par
\mbox{}\par
--- スケーリング: 安全な最短時間（転倒・滑りの両方）---\par
Earth: s= 1:  1.62 s, s= 2:  2.29 s, s= 5:  3.61 s, s=10:  5.11 s\par
Mars : s= 1:  2.62 s, s= 2:  3.71 s, s= 5:  5.87 s, s=10:  8.30 s\par
Moon : s= 1:  3.97 s, s= 2:  5.61 s, s= 5:  8.88 s, s=10: 12.55 s\par
T\_min / √(s · g\_E/g) = 1.616 〜 1.616（一定ならフルード則が厳密に成立）\par
\mbox{}\par
=== Part B: 腕の軌道最適化（転倒・滑り・関節速度/加速度の制約下で最短時間）===\par
単純軌道（最小躍度）: 1.62 s → 最適化軌道: 1.03 s（36\% 短縮）\par
同じ 1.03 s で実行した場合の ZMP: 単純 [-0.72, +0.78] m → 最適化 [-0.35, +0.35] m（接地範囲 ±0.4）\par
月面: 単純軌道 3.97 s → 最適化軌道 2.01 s\par
\mbox{}\par
=== Part C: 腕の水平振り（90°、速く振って 3 s で戻す）による台車のヨー ===\par
ヨー慣性: 台車 6.7, 腕+荷物 12.3 kg m², 平均半径 r̄ = 0.271 m\par
  Earth: 摩擦トルク上限  154.7 N m → 振りが 0.85 s より速いと台車が滑る\par
  Mars : 摩擦トルク上限   58.7 N m → 振りが 1.38 s より速いと台車が滑る\par
  Moon : 摩擦トルク上限   25.6 N m → 振りが 2.09 s より速いと台車が滑る\par
  摩擦ゼロ（宇宙と同じ）: 1 往復の正味ヨー max |1.9e-04|°（関節 1 つでは回らない）\par
  Earth: 1 往復の正味ヨー  T\_fast=0.25s: -37.03°, T\_fast=0.47s:  -8.16°, T\_fast=0.87s:  +0.00°, T\_fast=1.61s:  +0.00°, T\_fast=3.00s:  +0.00°\par
  Mars : 1 往復の正味ヨー  T\_fast=0.25s: -48.97°, T\_fast=0.47s: -31.18°, T\_fast=0.87s:  -2.24°, T\_fast=1.61s:  +0.00°, T\_fast=3.00s:  +0.00°\par
  Moon : 1 往復の正味ヨー  T\_fast=0.25s: -53.95°, T\_fast=0.47s: -43.70°, T\_fast=0.87s: -18.56°, T\_fast=1.61s:  -0.67°, T\_fast=3.00s:  +0.00°\par
  例: 0.6 s で振り 3 s で戻す往復を 5 回 → 地球 -2.8°, 月 -177.0°\par
}

\section{Sim07（\texttt{data/sim07\_summary.txt}）}
{\footnotesize\raggedright\sloppy\setlength{\parindent}{0pt}\setlength{\parskip}{0pt}\emergencystretch=3em
=== Part A1: 前方を軸に腕で円錐を描く（右腕、反時計回り = 軸方向から見て）===\par
  半頂角   5°: 1 周あたり胴体回転   0.162°, 回転軸 [-0.971 -0.135 -0.198]\par
  半頂角  10°: 1 周あたり胴体回転   0.647°, 回転軸 [-0.967 -0.161 -0.197]\par
  半頂角  20°: 1 周あたり胴体回転   2.575°, 回転軸 [-0.96  -0.207 -0.191]\par
  半頂角  30°: 1 周あたり胴体回転   5.732°, 回転軸 [-0.954 -0.24  -0.181]\par
  半頂角  45°: 1 周あたり胴体回転  12.440°, 回転軸 [-0.954 -0.257 -0.154]\par
  半頂角  60°: 1 周あたり胴体回転  20.432°, 回転軸 [-0.968 -0.224 -0.114]\par
  半頂角  70°: 1 周あたり胴体回転  25.357°, 回転軸 [-0.982 -0.169 -0.08 ]\par
  小さい円錐での指数: 回転角 ∝ β\^{}2.00（理論: 2 = 幾何学的位相の 2 次効果）\par
  逆回り（半頂角 45°）: [11.867  3.203  1.917]°（順回り [-11.867  -3.203  -1.917]° と符号反転）\par
\mbox{}\par
=== Part A2: 円錐の軸の向き（半頂角 45°）===\par
  forward (ψ=0, α=0)        :  12.44°/周, 回転軸 [-0.95 -0.26 -0.15]（円錐の軸との角度 17°）\par
  right side (ψ=-90, α=0)   :  11.08°/周, 回転軸 [-0.23  0.94 -0.25]（円錐の軸との角度 20°）\par
  up (α=90): azimuth sweep  :  58.03°/周, 回転軸 [ 0.    0.15 -0.99]（円錐の軸との角度 8°）\par
  front-up (ψ=0, α=45)      :  39.69°/周, 回転軸 [-0.25 -0.06 -0.97]（円錐の軸との角度 31°）\par
\mbox{}\par
=== Part A3: 球面座標の特異点（極 = 真上）の近くを通る円錐（T = 2 s）===\par
  極との最小角距離  30.0°: 最大関節速度 ψ'     3.14 rad/s, α' 1.57 rad/s（腕先端の角速度は一定 1.57 rad/s）\par
  極との最小角距離  10.0°: 最大関節速度 ψ'     9.05 rad/s, α' 1.57 rad/s（腕先端の角速度は一定 1.57 rad/s）\par
  極との最小角距離   3.0°: 最大関節速度 ψ'    30.01 rad/s, α' 1.57 rad/s（腕先端の角速度は一定 1.57 rad/s）\par
  極との最小角距離   1.0°: 最大関節速度 ψ'    90.00 rad/s, α' 1.57 rad/s（腕先端の角速度は一定 1.57 rad/s）\par
  極との最小角距離   0.3°: 最大関節速度 ψ'   300.00 rad/s, α' 1.57 rad/s（腕先端の角速度は一定 1.57 rad/s）\par
\mbox{}\par
=== Part A4: 回転の非可換性 ===\par
  ヨー用円錐 1 周（約 58°）とピッチ用円錐 3 周（約 33°）\par
  ヨー→ピッチ: 回転ベクトル [  8.2   41.44 -64.11]°\par
  ピッチ→ヨー: 回転ベクトル [-22.4   33.65 -65.27]°\par
  2 つの最終姿勢の差: 29.36°（平面ではループの順序を入れ替えても結果は同じ）\par
\mbox{}\par
=== Part B1: 静的な定格荷重マップ（台車 80 kg, 接地面 0.8×0.6 m, 腕 7 kg・1.0 m）===\par
  腕を水平に方位    0°（前）: 最大荷重 41.5 kg\par
  腕を水平に方位   45°（斜め前）: 最大荷重 42.2 kg\par
  腕を水平に方位   90°（横）: 最大荷重 24.3 kg\par
  腕を水平に方位  180°（後）: 最大荷重 41.5 kg\par
\mbox{}\par
=== Part B2: 腕（荷物 10 kg）を水平に 180° 旋回（前 → 左 → 後）===\par
  制約: 転倒（余裕 0.05 m）, 滑り（μ=0.6）, 関節速度 6.0 rad/s, 加速度 40.0 rad/s², 仰角 ≥ -35°\par
  Earth: 水平のまま 1.83 s / 途中で 60° 持ち上げ 1.24 s / 途中で 35° 下げる 1.01 s / 最適化 0.88 s（仰角の最大 0°, 最小 -35°）\par
  Moon : 水平のまま 4.48 s / 途中で 60° 持ち上げ 3.04 s / 途中で 35° 下げる 2.48 s / 最適化 2.07 s（仰角の最大 0°, 最小 -35°）\par
}

\section{Sim08（\texttt{data/sim08\_summary.txt}）}
{\footnotesize\raggedright\sloppy\setlength{\parindent}{0pt}\setlength{\parskip}{0pt}\emergencystretch=3em
=== Part A: コッドマンのパラドックス（ひねらずに腕を動かして元の向きに戻す）===\par
  下→前→横→下（古典的なコッドマンの動作）       : 上腕のひねり  +90.00°  / 囲んだ立体角  +90.00°（ラジアンを度に換算）\par
  腕で半頂角 20° の円錐を描いて下ろす        : 上腕のひねり  +21.71°  / 囲んだ立体角  +21.71°（ラジアンを度に換算）\par
  腕で半頂角 45° の円錐を描いて下ろす        : 上腕のひねり +105.44°  / 囲んだ立体角 +105.44°（ラジアンを度に換算）\par
  腕で半頂角 90° の円錐を描いて下ろす        : 上腕のひねり +360.00°  / 囲んだ立体角 +360.00°（ラジアンを度に換算）\par
  2 軸ジンバルの肩: 関節角が元に戻れば姿勢も元に戻る（ひねり 0°）。代わりに、ひねりを独立に選べない。\par
\mbox{}\par
=== Part B: 腕組み（右手を左上腕へ, 目標 [0.2  0.15 0.  ] m, 許容誤差 1.5 cm）===\par
  [細身 (奥行 0.18 m)  ] 2 軸ジンバル（Sim07）      : 不可（誤差  23.8 cm）  水平内転  22.8°, ひねり +157.2°（負=内旋）, 肘 100.2°\par
  [細身 (奥行 0.18 m)  ] 球関節（人の可動域）          : 不可（誤差   5.9 cm）  水平内転  14.1°, ひねり  -70.0°（負=内旋）, 肘  93.4°\par
  [細身 (奥行 0.18 m)  ] 球関節＋上肢帯の前方突出        : 不可（誤差   1.7 cm）  水平内転  45.0°, ひねり  -70.0°（負=内旋）, 肘  88.8°, 前方突出 12.2°\par
  [細身 (奥行 0.18 m)  ] 球関節（可動域制限なし）        : 可能（誤差   0.0 cm）  水平内転  16.8°, ひねり -101.9°（負=内旋）, 肘  84.6°\par
  [標準 (0.22 m)     ] 2 軸ジンバル（Sim07）      : 不可（誤差  25.1 cm）  水平内転 -157.4°, ひねり  -22.6°（負=内旋）, 肘 116.6°\par
  [標準 (0.22 m)     ] 球関節（人の可動域）          : 不可（誤差   6.6 cm）  水平内転  10.9°, ひねり  -70.0°（負=内旋）, 肘  94.9°\par
  [標準 (0.22 m)     ] 球関節＋上肢帯の前方突出        : 不可（誤差   2.6 cm）  水平内転  45.0°, ひねり  -70.0°（負=内旋）, 肘  90.6°, 前方突出 15.8°\par
  [標準 (0.22 m)     ] 球関節（可動域制限なし）        : 可能（誤差   0.0 cm）  水平内転  18.2°, ひねり +237.3°（負=内旋）, 肘  84.6°\par
  [厚い胸板・装甲 (0.26 m)] 2 軸ジンバル（Sim07）      : 不可（誤差  26.1 cm）  水平内転   0.0°, ひねり  -19.9°（負=内旋）, 肘 115.4°\par
  [厚い胸板・装甲 (0.26 m)] 球関節（人の可動域）          : 不可（誤差   7.1 cm）  水平内転   8.7°, ひねり  -70.0°（負=内旋）, 肘  96.1°\par
  [厚い胸板・装甲 (0.26 m)] 球関節＋上肢帯の前方突出        : 不可（誤差   3.6 cm）  水平内転  45.0°, ひねり  -70.0°（負=内旋）, 肘  92.5°, 前方突出 19.9°\par
  [厚い胸板・装甲 (0.26 m)] 球関節（可動域制限なし）        : 可能（誤差   1.4 cm）  水平内転  13.9°, ひねり +250.5°（負=内旋）, 肘  87.4°\par
\mbox{}\par
=== Part B2: 腕組みに最低限必要な内旋（ひねり）の大きさ（水平内転 ≤ 45°, 前方突出 ≤ 25° の下で）===\par
  参考: 人の内旋の可動域はおよそ 70〜90°（姿勢と測り方に依存）\par
  [細身 (奥行 0.18 m)  ] 球関節のみ    : 必要な内旋  86.0°（肘  87.1°）\par
  [細身 (奥行 0.18 m)  ] 球関節＋上肢帯: 必要な内旋  70.9°（肘  88.5°, 前方突出 12.2°）\par
  [標準 (0.22 m)     ] 球関節のみ    : 必要な内旋  92.7°（肘  87.3°）\par
  [標準 (0.22 m)     ] 球関節＋上肢帯: 必要な内旋  73.0°（肘  89.5°, 前方突出 15.8°）\par
  [厚い胸板・装甲 (0.26 m)] 球関節のみ    : 必要な内旋 109.5°（肘  87.4°）\par
  [厚い胸板・装甲 (0.26 m)] 球関節＋上肢帯: 必要な内旋  75.3°（肘  90.8°, 前方突出 19.9°）\par
}

\section{Sim09（\texttt{data/sim09\_summary.txt}）}
{\footnotesize\raggedright\sloppy\setlength{\parindent}{0pt}\setlength{\parskip}{0pt}\emergencystretch=3em
=== Part A1: 宇宙（平面）— 1 周あたりに稼げる回転の上限（総質量 76 kg で比較）===\par
  arm + leg (human)     : 総質量 76 kg → 上限  9.57°/周 → 90° に 10 周\par
  arm + arm (no legs)   : 総質量 76 kg → 上限  5.71°/周 → 90° に 16 周\par
  arm only (no legs)    : 総質量 76 kg → 関節 1 つなので正味の回転は 0°（往復で元に戻る）\par
  arm + heavy leg (x2)  : 総質量 76 kg → 上限 12.87°/周 → 90° に 7 周\par
\mbox{}\par
=== Part A2: 宇宙（3 次元）— 円錐運動 1 周あたりの胴体の回転 ===\par
  腕（肩）: 前方軸, 半頂角 45°              :  14.98°/周, 回転軸 [-0.97 -0.21 -0.11]\par
  腕（肩）: 前下方 30° 軸, 半頂角 30°        :  24.78°/周, 回転軸 [-0.19 -0.04  0.98]\par
  脚（股関節）: 前下方 30° 軸, 半頂角 30°      :  48.23°/周, 回転軸 [-0.23 -0.13  0.96]\par
  脚（股関節）: 前下方 30° 軸, 半頂角 45°      :  96.78°/周, 回転軸 [-0.21 -0.21  0.95]\par
  ※ 股関節の可動域は肩より狭い（屈曲 ≈120°, 伸展 ≈20°, 外転 ≈45°）。同じ軸・半頂角（前下方 30°, 30°）で比べると脚は腕の 1.9 倍回る\par
\mbox{}\par
=== Part B: 地上 — 荷物運び（Sim06 と同じ腕の動作, 荷物 10 kg）の転ばない最短時間 ===\par
  支持範囲: 戦車 ±0.4 m, 二本脚（足裏） ±0.125 m / 安全余裕 0.02 m\par
  tank  : 静的な ZMP の範囲 [-0.108, +0.133] m, 単純軌道の最短 1.52 s\par
  biped : 静的な ZMP の範囲 [-0.108, +0.133] m, 単純軌道の最短 inf s\par
  tank      : 腕の軌道を最適化 → 最短 0.98 s\par
  tank+hip  : 腕の軌道を最適化＋上体の前後傾 → 最短 0.64 s（上体の傾き -32°〜+29°, 負=後ろに反る）\par
  biped     : 腕の軌道を最適化 → 解なし（静止状態でも転倒）\par
  biped+hip : 腕の軌道を最適化＋上体の前後傾 → 最短 0.69 s（上体の傾き -40°〜+28°, 負=後ろに反る）\par
}

\section{Sim10（\texttt{data/sim10\_summary.txt}）}
{\footnotesize\raggedright\sloppy\setlength{\parindent}{0pt}\setlength{\parskip}{0pt}\emergencystretch=3em
=== Part A: 戦車型（履帯接地長 0.8 m, 前端の高さ 0.25 m, μ = 0.6）===\par
  arm stowed back-up (COM rear) : 重心 x -0.023 m, 高さ 0.407 m → 力の限界 0.26 m, 越えられる段差 0.25 m（20 cm の段差に必要な摩擦係数 0.16）\par
  arm down (neutral)            : 重心 x +0.000 m, 高さ 0.345 m → 力の限界 0.30 m, 越えられる段差 0.25 m（20 cm の段差に必要な摩擦係数 0.23）\par
  arm forward (COM front)       : 重心 x +0.037 m, 高さ 0.382 m → 力の限界 0.33 m, 越えられる段差 0.25 m（20 cm の段差に必要な摩擦係数 0.24）\par
\mbox{}\par
=== Part B: 二本脚（股関節高さ 0.9 m, 大腿・下腿 0.45/0.45 m, 膝・股関節トルク上限 250/250 N m）===\par
  可動域の限界（足を段の上に置けるか）: 0.63 m\par
  トルクの限界: 上体を立てたまま 0.18 m / 上体の前傾で重心を前へ移す 0.60 m\par
    段差 0.15 m で上体を立てたままの膝トルク    227 N m\par
    段差 0.20 m で上体を立てたままの膝トルク    258 N m\par
    段差 0.30 m で上体を立てたままの膝トルク    306 N m\par
    段差 0.40 m で上体を立てたままの膝トルク    341 N m\par
  → 越えられる段差 0.60 m（上体の前傾あり。ゆっくり上る場合で、蹴り出しや勢いは使わない）\par
\mbox{}\par
=== Part C: 相似拡大（長さ s 倍, 質量 s³ 倍, トルク上限 ∝ s³）での越えられる段差 / 機体の大きさ ===\par
  s =  0.5（全高 約  0.9 m）: 戦車   0.12 m（大きさ比 0.25）, 二本脚   0.32 m（大きさ比 0.63, 可動域 0.63 / トルク 0.90）\par
  s =  1.0（全高 約  1.8 m）: 戦車   0.25 m（大きさ比 0.25）, 二本脚   0.60 m（大きさ比 0.60, 可動域 0.63 / トルク 0.60）\par
  s =  2.0（全高 約  3.6 m）: 戦車   0.50 m（大きさ比 0.25）, 二本脚   0.37 m（大きさ比 0.18, 可動域 0.63 / トルク 0.18）\par
  s =  5.0（全高 約  9.0 m）: 戦車   1.25 m（大きさ比 0.25）, 二本脚   0.12 m（大きさ比 0.02, 可動域 0.63 / トルク 0.02）\par
  s = 10.0（全高 約 18.0 m）: 戦車   2.50 m（大きさ比 0.25）, 二本脚   0.05 m（大きさ比 0.00, 可動域 0.63 / トルク 0.00）\par
}

\section{詳細版の追加解析（\texttt{data/detail\_summary.txt}）}
{\footnotesize\raggedright\sloppy\setlength{\parindent}{0pt}\setlength{\parskip}{0pt}\emergencystretch=3em
[D1] 閉じた形と数値計算の最大差 1.78e-15（A の範囲 5.776〜7.613, Γ の範囲 0.0283〜0.1421）\par
[D2] 腕 1 本で得られる最大の首振り: 32.1°（下 → 上の 180° で -16.0°）\par
[D3] 1 周あたりの上界（総質量 76 kg）: 腕4kg・脚6kg → 6.41°, 腕4kg・脚12kg → 9.60°, 腕4kg・脚24kg → 12.91°, 腕8kg・脚12kg → 16.90°\par
[D4] RW 0.5 kg: 90° に必要なホイール回転 1598 回, 10 s で旋回するときの最大 17982 rpm\par
[D4] RW 2.0 kg: 90° に必要なホイール回転 178 回, 10 s で旋回するときの最大 2003 rpm\par
[D4] RW 8.0 kg: 90° に必要なホイール回転 20 回, 10 s で旋回するときの最大 226 rpm\par
[D5] 股関節トルク ∝ f\^{}2.00, 90° 旋回の正仕事 ∝ f\^{}2.00（0.1 Hz で 19.9 J, 2 Hz で 7946 J）\par
[D6] 学習した位相速度 u(e): e=+90° で +1.00, e=−90° で -1.00, e=0 付近で -0.014\par
[D7] T=20 s: 推進剤 9.37 kg, スラスタの噴射は最初と最後の各 0.5 s 程度, AMBAC は全区間で最大周波数（平均 f/fmax 1.00）\par
[D7] T=60 s: 推進剤 1.97 kg, スラスタの噴射は最初と最後の各 0.5 s 程度, AMBAC は全区間で最大周波数（平均 f/fmax 1.00）\par
[D7] T=110 s: 推進剤 0.32 kg, スラスタの噴射は最初と最後の各 0.6 s 程度, AMBAC は全区間で最大周波数（平均 f/fmax 1.00）\par
[D8] T = 1.616√(s·g\_E/g) で、4 条件の規格化 ZMP が一致（図）\par
[D9] 月面: μ=0.2 で 0.3 s の振り → -50.5°/往復, μ=1.0 で → -48.6°/往復\par
[D10] 半頂角 30° の円錐: 1 周あたり 5.1〜38.2°（最大は軸 ψ=90°, α=-60°）\par
[D11] 胸の奥行 → 必要な内旋（球関節のみ / ＋上肢帯, 前方突出）: 0.16 m → 83° / 70° (11°), 0.18 m → 86° / 71° (12°), 0.20 m → 89° / 72° (14°), 0.22 m → 93° / 73° (16°), 0.24 m → 97° / 74° (18°), 0.26 m → 110° / 75° (20°), 0.28 m → nan° / 77° (22°)\par
[D12] 脚の質量の割合 2\% → 2.2°/周, 40\% → 14.3°/周（人間型 16\% 付近 ≈ 9.6°）\par
[D13] 戦車型の力の限界（腕を前）: μ=0.1 → 0.05 m, μ=0.15 → 0.08 m, μ=0.2 → 0.33 m, μ=0.25 → 0.33 m, μ=0.3 → 0.33 m, μ=0.4 → 0.33 m, μ=0.6 → 0.33 m, μ=1.0 → 0.33 m\par
[D13] 二本脚の限界（上体の前傾あり）: 100 N m → 0.11 m, 175 N m → 0.31 m, 250 N m → 0.60 m, 325 N m → 0.90 m, 400 N m → 0.90 m\par
[D14] 6 種類のループの 1 周あたりの回転: +6.70°, +1.44°, +1.25°, +0.57°, +0.71°, +0.67°\par
[D14] 相対誤差 1\% あたりの回転の変化 [°]: arm mass: +0.015, +0.003, +0.002, +0.001, +0.002, +0.001; leg mass: +0.003, +0.000, +0.000, +0.000, +0.000, +0.000; torso inertia: -0.079, -0.016, -0.018, -0.008, -0.003, -0.011\par
[D14] 計測誤差 0.1°・各 10 周で推定した場合の相対標準偏差: 腕 29.9\%, 脚 223.8\%, 胴体の慣性 4.1\%（条件数 603.8）\par
[D14] 周回中の姿勢の波形（各ループ 36 点）も使う場合の相対標準偏差: 腕 3.20\%, 脚 10.07\%, 胴体の慣性 0.94\%（条件数 316.0）\par
}
